\documentclass[preprint,12pt]{elsarticle}
\usepackage{amsmath,amssymb}
\usepackage{graphicx}
\usepackage{hyperref}
\usepackage{booktabs}

\begin{document}

\begin{frontmatter}
\title{Volumetric neutral reactions in OpenEdge: charge exchange and dissociation for SOLPS--OpenEdge coupled simulations}

\author[ornl]{A.~Diaw\corref{cor1}}
\ead{diawa@ornl.gov}
\cortext[cor1]{Corresponding author}
\address[ornl]{Oak Ridge National Laboratory, Oak Ridge, TN 37831, USA}

\begin{abstract}
We describe the implementation of volumetric neutral transport reactions in OpenEdge,
a kinetic impurity and neutral particle transport code built on the SPARTA framework.
Charge exchange, electron-impact dissociation, and ionization reactions are added to the
existing ADAS-based chemistry module, enabling OpenEdge to replace EIRENE's role in
coupled SOLPS--OpenEdge simulations. Rate coefficients use ADAS ADF11 tables for
ionization, recombination, and charge exchange, and Janev--Langer--Evans--Post polynomial
fits for molecular dissociation. The implementation supports GPU acceleration via Kokkos.
\end{abstract}

\begin{keyword}
neutral transport \sep charge exchange \sep dissociation \sep DSMC \sep plasma--material interaction
\end{keyword}
\end{frontmatter}

%% ============================================================
\section{Introduction}
\label{sec:intro}

Neutral transport in the scrape-off layer (SOL) of magnetic fusion devices plays a critical
role in plasma--material interactions, fuel recycling, and impurity generation. EIRENE~\cite{Reiter2005}
is the standard Monte Carlo neutral transport code used in the coupled B2.5--EIRENE (SOLPS)
framework. However, EIRENE's straight-line free-flight model, file-based coupling interface,
and species count limitations motivate the development of an alternative within OpenEdge.

OpenEdge is a kinetic particle transport code built as a package on top of
SPARTA~\cite{Gallis2014}, providing Boris and guiding-center-approximation (GCA) pushers,
background plasma interpolation, and surface interaction models. By adding volumetric
neutral reactions directly to OpenEdge, we can:
\begin{itemize}
  \item Eliminate the file-based EIRENE coupling loop,
  \item Leverage the same MPI + GPU (Kokkos) parallel infrastructure used for impurity transport,
  \item Use point-query plasma interpolation at particle positions rather than cell-averaged fields.
\end{itemize}

%% ============================================================
\section{Reaction framework}
\label{sec:framework}

The existing \texttt{fix chem/adas} command in OpenEdge handles ionization and recombination
using ADAS ADF11 rate coefficient tables. We extend this framework with two additions:
(1) charge exchange using CCD rate tables, and (2) molecular dissociation using
Janev polynomial fits.

\subsection{Rate coefficient evaluation}

For ADAS-based reactions (ionization, recombination, charge exchange), the rate coefficient
$\langle\sigma v\rangle(T_e, n_e)$ is obtained by bilinear interpolation on a 2D grid in
$(\log_{10} T_e, \log_{10} n_e)$ space. The rate tables are stored in HDF5 files indexed by
charge state $q$, electron temperature, and electron density.

For Janev-style reactions (dissociation), the rate coefficient is evaluated from a 9-term
polynomial fit:
\begin{equation}
  \ln\langle\sigma v\rangle = \sum_{n=0}^{8} b_n (\ln T_e)^n
  \label{eq:janev}
\end{equation}
where $T_e$ is in eV and $\langle\sigma v\rangle$ is in cm$^3$/s. Coefficients are taken
from the HYDHEL database~\cite{Reiter1987,Janev1987}.

\subsection{Poisson competing-channel selection}

At each chemistry timestep $\Delta t_{\text{chem}}$, each particle is subjected to all
available reaction channels simultaneously. For each channel $i$, the Poisson rate parameter is:
\begin{equation}
  \lambda_i = \langle\sigma v\rangle_i \cdot n_e \cdot \Delta t_{\text{chem}}
\end{equation}
The probability of at least one reaction occurring is $P = 1 - e^{-\lambda_{\text{total}}}$,
and the channel is selected proportionally to $\lambda_i / \lambda_{\text{total}}$.

%% ============================================================
\section{Charge exchange}
\label{sec:cx}

The charge exchange reaction for an impurity ion $A^{z+}$ with a neutral hydrogen donor:
\begin{equation}
  A^{z+} + \text{H} \rightarrow A^{(z-1)+} + \text{H}^+
\end{equation}
reduces the charge state by one, analogous to recombination. The CCD rate coefficients
from ADAS are stored in the same HDF5 format as ionization (SCD) and recombination (ACD)
tables, and are read by the same bilinear interpolation infrastructure.

For hydrogen charge exchange ($\text{D}^+ + \text{D} \rightarrow \text{D} + \text{D}^+$),
the rate is $\sim 10^{-9}$~cm$^3$/s at $T_e = 10$--100~eV, dominating radiative
recombination ($\sim 10^{-13}$~cm$^3$/s) by several orders of magnitude.

%% ============================================================
\section{Electron-impact dissociation}
\label{sec:dissociation}

The primary molecular deuterium dissociation channel is:
\begin{equation}
  e + \text{D}_2 \rightarrow 2\text{D} + e
  \label{eq:dissoc}
\end{equation}
corresponding to HYDHEL reaction H.2~2.2.5. This is an electron-impact process where
the D$_2$ molecule dissociates via excitation to the repulsive $b\,^3\Sigma_u^+$ state.

In the particle framework, dissociation converts one D$_2$ particle into two D particles.
The first product replaces the reactant; the second is created at the same position with
the same velocity (to be refined with thermal velocity sampling in future work).

The rate coefficient (Eq.~\ref{eq:janev}) with HYDHEL coefficients gives
$\langle\sigma v\rangle \approx 4 \times 10^{-9}$~cm$^3$/s at $T_e = 20$~eV,
rising steeply from a threshold near $\sim 1$~eV.

%% ============================================================
\section{Implementation}
\label{sec:implementation}

The reaction framework is implemented in the \texttt{fix chem/adas} SPARTA fix class.
Key design choices:
\begin{itemize}
  \item \textbf{Rate styles:} \texttt{A} (ADAS table lookup) and \texttt{J} (Janev polynomial)
    are dispatched per reaction in the competing-channel loop.
  \item \textbf{Particle creation:} Dissociation products are deferred to avoid invalidating
    the particle array during iteration, then created after the per-cell loop completes.
  \item \textbf{GPU acceleration:} The Kokkos kernel (\texttt{fix chem/adas/kk}) runs rate
    interpolation and Poisson draws entirely on device. Charge exchange uses the same
    device-side bilinear interpolation as ionization and recombination.
  \item \textbf{Backward compatibility:} HDF5 files without CX data are handled gracefully
    (existence check before reading).
\end{itemize}

%% ============================================================
\section{SOLPS--OpenEdge coupling}
\label{sec:coupling}

With volumetric neutral reactions, OpenEdge can replace EIRENE in the SOLPS coupling loop.
The plasma background ($T_e$, $n_e$, flows) is provided by B2.5 via HDF5 files and
interpolated at particle positions. OpenEdge returns ionization source terms (particle
and momentum) to B2.5 for the next fluid iteration.

This eliminates the EIRENE file-based interface and its associated species count limitations
(INDPRO errors), while enabling GPU-parallel neutral transport on the same mesh as impurity
ion transport.

%% ============================================================
\section{Conclusions}
\label{sec:conclusions}

We have implemented charge exchange and electron-impact dissociation in OpenEdge's ADAS
chemistry framework, completing the core set of volumetric reactions needed for neutral
transport in fusion edge plasmas. The implementation reuses the existing rate interpolation
and Poisson channel selection infrastructure, with Janev polynomial fits added for molecular
reactions. Future work includes velocity re-sampling from the background ion Maxwellian
after CX and dissociation events, and surface recycling models.

\section*{Acknowledgments}
This work was supported by the U.S. Department of Energy.

\bibliographystyle{elsarticle-num}
\bibliography{references}

\end{document}
